\subsection{本章目标与边界：导出骨架，而非现成物理方程}\label{subsec:physical-spacetime-skeleton-goal-boundary}

本章不引入新的时空公理，不增设新的基本对象，也不在统一上位逻辑之外附加一个外部的物理背景。本章所做的，只是把前文已经建立的观察者、事件、区域、状态更新、局部 forcing、\(\mathsf{NULL}\)、残差、谱接口与资源代价放到同一条解释链上重新阅读，并说明：当这些结构被联立起来时，体系内部已经出现了一种可被物理化理解的时空骨架。这里所谓“骨架”，并不是一个现成的几何时空，而只是使观察者定位、事件排序、局部比较、时间投影、共享支撑与资源代价同时可表达的最小结构。

这一判断并不依赖新的假设。第 \ref{sec:logic-expansion-chain} 节已经给出了观察者索引状态、事件—区域系统、因果相容更新与共享 forcing 的抽象骨架，尤其见 \S\ref{subsec:logic-expansion-chain-observer-spacetime} 与 \S\ref{subsec:logic-expansion-chain-instantiation-interface}；第 \ref{sec:recursive-addressing} 节把局部截面、\(\mathsf{NULL}\) 与延迟分类组织为对象层可断言性的语义，并以命题 \ref{prop:observer-indexed-delayed-classification-not-retrocausal} 排除了“延迟分类回写过去”的误读，以定理 \ref{thm:prefix-site-cech-null-gerbe} 给出局部存在与全局失败之间的 \v{C}ech 型粘接障碍；第 \ref{sec:pom} 节把“提升—演化—投影—证书”的生成链压缩到同一接口内，并在 \S\ref{subsec:pom-time-space-diagram}、\S\ref{subsec:pom-fold-compression} 与 \S\ref{subsec:pom-conditional-expectation} 中给出时间—空间交换图、折叠压缩率与条件期望接口；第 \ref{sec:circle-dimension-phase-gate} 节则把这一局部截面—粘接障碍口径继续推进到相位与精度接口之中。第 \ref{sec:fold-residual-time} 节进一步将残差展开、序列层可逆化与最小在线延迟纳入同一时间化口径，见 \S\ref{subsec:residual-unfolding-time-depth} 与 \S\ref{subsubsec:frt-min-online-delay}。因此，本章不是另起炉灶，而是对前文结果作一次显式收束：前文已经给出了足以承载物理时空解释的全部核心部件，本章只把这些部件组合成一条清楚的主线。

因此，本章的立场与“先设定一个时空容器，再把观察者和测量放入其中”的通常写法不同。在这里，观察者首先不是特权主体，而是状态的索引纤维；事件首先不是几何点，而是状态更新与命题定位的发生位置；区域首先不是绝对容器，而是局部 forcing、局部截面与共享支撑得以成立的覆盖片段；时间首先不是独立背景，而是状态精化、残差外置与事件定位共同诱导出的顺序投影。换言之，在本体系中，时空不是论证的前提，而是统一上位逻辑在观察者化、局部化、因果化与资源化之后的一种解释性收缩。

这也同时规定了本章的边界。本章不直接推出 Minkowski 度规，不直接推出 Lorentz 对称，不直接推出 Einstein 型场方程，也不直接推出任何具体测量公设。凡此种种，都属于更强的实例化层结论；它们只有在某个具体数学对象的实现同时保持 forcing 结构、局部性、因果相容、失败模式、时间—空间接口与资源代价之后，才可能由本章所抽取的骨架进一步收缩出来。因此，本章的目标不是给出一套现成物理学，而是说明：统一上位逻辑已经给出了一个足以承载物理时空解释的最小前体。

本章的推理顺序也因此是明确的。首先，从观察者纤维、事件—区域系统与因果相容更新中抽取最小的局域因果骨架；其次，从状态精化链、残差外化与最小在线延迟中抽取时间的投影性；再次，从共享 forcing、局部截面与共同区域支撑中说明空间首先表现为局部可比较性，而不是绝对容器；最后，再借助时间—空间交换图、primitive 轨道谱、资源代价与 \v{C}ech 型粘接障碍，说明何种实例化有可能进一步产生度规、规范或曲率的读法。于是，本章并不添加新的逻辑内容，而是在不改变前文结论的前提下，把前文已经分散出现的结果收束为一条通向物理时空的解释链。

概括地说，本章所要确立的并不是“物理时空已被预设”，而是相反的命题：统一上位逻辑先给出观察者化的可断言结构、局部支撑结构、因果结构、谱接口与资源接口，而我们通常称之为“物理时空”的东西，只是在某些满足条件的实例化中，由这些结构共同涌现出来的稳定读出。其逻辑主链可以概括为：状态—forcing 结构，经由观察者纤维化而获得局部视角；经由事件—区域—覆盖而获得局部组织；经由因果相容更新而获得先后关系；经由残差外置与最小在线延迟而获得时间投影；经由共享 forcing 与局部截面而获得共同支撑与区域共识；再经由谱接口与资源代价而获得进一步的几何候选。到这一步，所谓“物理时空骨架”已经不是外加假设，而是前文结构联立后的自然读法。

后续各小节将依次展开这条主链：先说明观察者为何只是一种状态纤维而非特权主体，再说明事件—区域—覆盖何以构成最小时空骨架，继而讨论时间的投影性、空间的共同支撑起源，以及谱与资源结构如何把这一骨架推进到可与具体物理实例化对接的层面。
